\documentclass[10pt,a4paper]{amsart}
\usepackage[margin=19mm]{geometry}
\usepackage{amsmath,amssymb,mathtools}
\usepackage[scr=rsfs,cal=euler]{mathalfa}
\usepackage{xfrac}
\usepackage[unicode,hidelinks,bookmarks=false]{hyperref}
\newcommand{\ie}{i.e.\ }
\newcommand{\cf}{cf.\ }
\newcommand{\resp}{respectively\ }
\newcommand{\Subsection}{\subsection}
\providecommand{\nobreakdash}{\nobreak\hbox{-}}
\setlength{\emergencystretch}{2em}
\multlinegap0pt
\usepackage{bm}
\usepackage{bbm}
\usepackage{dsfont}
\usepackage{xspace}
\usepackage{cancel}
\usepackage{mathtools}
\usepackage[shortlabels]{enumitem}
\usepackage{amsthm}
\makeatletter
\@ifundefined{theorem}{\newtheorem{theorem}{Theorem}}{}
\makeatother
\title[FusedANN: Convexified Hybrid ANN via Attribute-Vector Fusion]{FusedANN: Convexified Hybrid ANN via Attribute-Vector Fusion}
\author{Alireza Heidari, Wei Zhang, Ying Xiong}
\date{Source: arXiv:2509.19767v2 (CC BY 4.0)}
\begin{document}
\maketitle

\noindent\emph{Source and licence.} FusedANN: Convexified Hybrid ANN via Attribute-Vector Fusion. Version arXiv:2509.19767v2 (CC BY 4.0), distributed under the Creative Commons Attribution 4.0 International licence, as recorded in the arXiv licence metadata for this exact version and shown on the abstract page https://arxiv.org/abs/2509.19767v2. Full rights notice: licenses/arxiv-2509.19767-CC-BY-4.0.txt.\par\smallskip

\noindent\textit{Editorial scope.} Counted unit: the Theorem whose source label is \texttt{theo:psi\_property} in the article (source0.tex lines 1126--1147, source label \texttt{theo:psi\_property}), delivered together with the article's own complete original proof of it (source0.tex lines 1149--1274, one \texttt{proof} environment of 126 lines). No mathematics is written, repaired or re-derived: statement and proof are the article's own text line for line. Required context retained uncounted: none. Every definition, display and label of those lines is retained unchanged. Omitted from this package and not redistributed: 1--1125, 1148--1148, 1275--3039. Nothing inside a retained excerpt is omitted or altered: every retained body is delivered exactly as the article's lines are written, so no omission receipt of any kind accompanies this package. Citations: the retained excerpts carry no citation command at all, so no bibliography entry of the article is transcribed and no reference is redistributed. Reference adaptation: none was needed and none was made; every reference inside the delivered unit resolves inside it, so no unresolved reference or citation is produced. Printed slips kept literal: no printed word, symbol, hypothesis or number of the counted statement or of its proof was altered. Layout and class: the article's own class is \texttt{article} with packages including neurips\_2025, inputenc, fontenc, hyperref, float, amsmath, amssymb, amsfonts, amsthm, nicefrac, microtype, algorithm, algorithmic, tikz; the wrapper is the lane's pinned \texttt{amsart} editorial wrapper at 10pt on A4, which restates the theorem-like environments the retained text needs and adds the article's own macro definitions that the retained text needs (none), with extra wrapper packages (bm, bbm, dsfont, xspace, cancel, mathtools, enumitem). The delivered numbering is the unit's own. Assets: no retained span contains a figure environment, a graphics-inclusion command, a PDF-page inclusion, a table or a TikZ picture; no image file is copied into this package, the package declares no asset dependency and no image file is redistributed with it. Licence: version arXiv:2509.19767v2 of this article is distributed under the Creative Commons Attribution 4.0 International licence, as recorded in the arXiv licence metadata for this exact version and shown on the abstract page https://arxiv.org/abs/2509.19767v2. A full rights notice is delivered at licenses/arxiv-2509.19767-CC-BY-4.0.txt. Characters: every retained excerpt is pure ASCII, so the delivered engine, XeLaTeX with its default Latin Modern text font, reports no missing character.
\begin{theorem}[Properties of $\Psi$ Transformation]
\label{theo:psi_property}
Let $\mathcal{D}$ be a record set with content vectors in $\mathbb{R}^d$ and attribute vectors in $\mathbb{R}^m$ where $m < d$ and $m \mid d$. For records $o_i, o_j \in \mathcal{D}$, let $v'_i = \Psi(v(o_i), f(o_i), \alpha, \beta)$ and $v'_j = \Psi(v(o_j), f(o_j), \alpha, \beta)$ be their transformed vectors under:
\begin{equation}
\Psi(v, f, \alpha, \beta) = \left[\frac{v^{(1)} - \alpha f}{\beta}, \ldots, \frac{v^{(d/m)} - \alpha f}{\beta}\right]
\end{equation}
where $\alpha > 1$ and $\beta > 0$ are scaling parameters. Then:

\begin{enumerate}[leftmargin=2em]
\item \textbf{Order Preservation for Same Attributes:} For any query $q$ with content vector $v(q)$ and attribute vector $f(q)$ where $f(q) = f(o_i) = f(o_j)$, if $\rho(v(o_i), v(q)) < \rho(v(o_j), v(q))$ in the original space, then $\rho(v'_i, v'_q) < \rho(v'_j, v'_q)$ in the transformed space, where $v'_q = \Psi(v(q), f(q), \alpha, \beta)$.

\item \textbf{Distance Preservation:} If $\beta = 1$, then $\rho(v'_i, v'_j) = \rho(v(o_i), v(o_j))$ for all $o_i, o_j$ with identical attributes.

\item \textbf{Attribute Separation:} For records $o_i, o_j$ with different attributes $f(o_i) \neq f(o_j)$, the distance $\rho(v'_i, v'j)$ increases as $\alpha$ increases, with a lower bound:
\begin{equation}
\rho(v'_i, v'_j) \geq \frac{1}{\beta}\sqrt{\rho^2(v(o_i), v(o_j)) + \alpha^2 \cdot \frac{d}{m} \cdot \rho^2(f(o_i), f(o_j)) - 2\alpha \cdot C{ij}}
\end{equation}
where $C_{ij} = |\sum_{l=1}^{d/m} \langle v^{(l)}(o_i) - v^{(l)}(o_j), f(o_i) - f(o_j) \rangle|$.

\item \textbf{Attribute Distance Order Preservation:} For records with identical content vectors but different attributes ($v(o_i) = v(o_j) = v(o_k) = v(o_l)$ but $f(o_i) \neq f(o_j)$ and $f(o_k) \neq f(o_l)$), if $\rho(f(o_i), f(o_j)) < \rho(f(o_k), f(o_l))$, then $\rho(v'_i, v'_j) < \rho(v'_k, v'_l)$.
\end{enumerate}
\end{theorem}

\begin{proof}
\textbf{Part 1: Order Preservation for Same Attributes.}
Consider records $o_i$ and $o_j$ with identical attributes $f(o_i) = f(o_j) = f$. Their transformed vectors are:
\begin{align}
v'_i &= \Psi(v(o_i), f, \alpha, \beta) = \left[\frac{v^{(1)}(o_i) - \alpha f}{\beta}, \ldots, \frac{v^{(d/m)}(o_i) - \alpha f}{\beta}\right] \\
v'_j &= \Psi(v(o_j), f, \alpha, \beta) = \left[\frac{v^{(1)}(o_j) - \alpha f}{\beta}, \ldots, \frac{v^{(d/m)}(o_j) - \alpha f}{\beta}\right]
\end{align}

Let us compute the squared Euclidean distance between these transformed vectors:
\begin{align}
\rho^2(v'i, v'j) &= \sum_{l=1}^{d/m} \sum_{h=1}^{m} \left(\frac{v^{(l)}(o_i)[h] - \alpha f[h]}{\beta} - \frac{v^{(l)}(o_j)[h] - \alpha f[h]}{\beta}\right)^2 \\
&= \sum_{l=1}^{d/m} \sum_{h=1}^{m} \left(\frac{v^{(l)}(o_i)[h] - v^{(l)}(o_j)[h]}{\beta}\right)^2 \\
&= \frac{1}{\beta^2}\sum_{l=1}^{d/m} \sum_{h=1}^{m} \left(v^{(l)}(o_i)[h] - v^{(l)}(o_j)[h]\right)^2 \\
&= \frac{1}{\beta^2}\sum_{p=0}^{d-1} \left(v(o_i)[p] - v(o_j)[p]\right)^2 \\
&= \frac{1}{\beta^2}\rho^2(v(o_i), v(o_j))
\end{align}

Taking the square root of both sides:
\begin{equation}
\rho(v'_i, v'_j) = \frac{1}{\beta}\rho(v(o_i), v(o_j))
\end{equation}

Now consider a query $q$ with $f(q) = f$. The transformed query vector is $v'_q = \Psi(v(q), f, \alpha, \beta)$. By the same derivation:
\begin{align}
\rho(v'_i, v'_q) &= \frac{1}{\beta}\rho(v(o_i), v(q)) \\
\rho(v'_j, v'_q) &= \frac{1}{\beta}\rho(v(o_j), v(q))
\end{align}

Since $\beta > 0$, the scaling factor $\frac{1}{\beta}$ preserves the inequality. Therefore:
\begin{equation}
\rho(v(o_i), v(q)) < \rho(v(o_j), v(q)) \Rightarrow \rho(v'_i, v'_q) < \rho(v'_j, v'_q)
\end{equation}

This establishes that the order of k-nearest neighbors is preserved for records with identical attributes.

\textbf{Part 2: Distance Preservation.}
When $\beta = 1$, the equation derived in Part 1 simplifies to:
\begin{equation}
\rho(v'_i, v'_j) = \rho(v(o_i), v(o_j))
\end{equation}

Therefore, if $\beta = 1$, the distances between records with identical attributes are exactly preserved.

\textbf{Part 3: Attribute Separation.}
For records $o_i$ and $o_j$ with different attributes $f(o_i) \neq f(o_j)$, their transformed vectors are:
\begin{align}
v'_i &= \Psi(v(o_i), f(o_i), \alpha, \beta) = \left[\frac{v^{(1)}(o_i) - \alpha f(o_i)}{\beta}, \ldots, \frac{v^{(d/m)}(o_i) - \alpha f(o_i)}{\beta}\right] \\
v'_j &= \Psi(v(o_j), f(o_j), \alpha, \beta) = \left[\frac{v^{(1)}(o_j) - \alpha f(o_j)}{\beta}, \ldots, \frac{v^{(d/m)}(o_j) - \alpha f(o_j)}{\beta}\right]
\end{align}

The squared Euclidean distance between these transformed vectors is:
\begin{align}
\rho^2(v'i, v'j) &= \sum_{l=1}^{d/m} \sum_{h=1}^{m} \left(\frac{v^{(l)}(o_i)[h] - \alpha f(o_i)[h]}{\beta} - \frac{v^{(l)}(o_j)[h] - \alpha f(o_j)[h]}{\beta}\right)^2 \\
&= \frac{1}{\beta^2}\sum_{l=1}^{d/m} \sum_{h=1}^{m} \left(v^{(l)}(o_i)[h] - v^{(l)}(o_j)[h] - \alpha(f(o_i)[h] - f(o_j)[h])\right)^2
\end{align}

Expanding the squared term:
\begin{align}
\rho^2(v'i, v'j) &= \frac{1}{\beta^2}\sum_{l=1}^{d/m} \sum_{h=1}^{m} \Big[ (v^{(l)}(o_i)[h] - v^{(l)}(o_j)[h])^2 + \alpha^2(f(o_i)[h] - f(o_j)[h])^2 \\
&\quad - 2\alpha(v^{(l)}(o_i)[h] - v^{(l)}(o_j)[h])(f(o_i)[h] - f(o_j)[h]) \Big] \\
&= \frac{1}{\beta^2}\Big[\rho^2(v(o_i), v(o_j)) + \alpha^2 \sum_{l=1}^{d/m} \sum_{h=1}^{m} (f(o_i)[h] - f(o_j)[h])^2 \\
&\quad - 2\alpha \sum_{l=1}^{d/m} \sum_{h=1}^{m} (v^{(l)}(o_i)[h] - v^{(l)}(o_j)[h])(f(o_i)[h] - f(o_j)[h])\Big]
\end{align}

Note that:
\begin{equation}
\sum_{l=1}^{d/m} \sum_{h=1}^{m} (f(o_i)[h] - f(o_j)[h])^2 = \frac{d}{m} \cdot \rho^2(f(o_i), f(o_j))
\end{equation}

And for the cross-term:
\begin{equation}
\sum_{l=1}^{d/m} \sum_{h=1}^{m} (v^{(l)}(o_i)[h] - v^{(l)}(o_j)[h])(f(o_i)[h] - f(o_j)[h]) = \sum_{l=1}^{d/m} \langle v^{(l)}(o_i) - v^{(l)}(o_j), f(o_i) - f(o_j) \rangle
\end{equation}

Let $C_{ij} = |\sum_{l=1}^{d/m} \langle v^{(l)}(o_i) - v^{(l)}(o_j), f(o_i) - f(o_j) \rangle|$. The squared distance becomes:
\begin{align}
\rho^2(v'_i, v'j) &= \frac{1}{\beta^2}\Big[\rho^2(v(o_i), v(o_j)) + \alpha^2 \cdot \frac{d}{m} \cdot \rho^2(f(o_i), f(o_j)) - 2\alpha \cdot C{ij}\Big]
\end{align}

Taking the derivative with respect to $\alpha$:
\begin{align}
\frac{\partial}{\partial \alpha}\rho^2(v'i, v'j) &= \frac{1}{\beta^2}\Big[2\alpha \cdot \frac{d}{m} \cdot \rho^2(f(o_i), f(o_j)) - 2C{ij}\Big]\
&= \frac{2}{\beta^2}\Big[\alpha \cdot \frac{d}{m} \cdot \rho^2(f(o_i), f(o_j)) - C{ij}\Big]
\end{align}

Since $\alpha > 1$ and $\frac{d}{m} \cdot \rho^2(f(o_i), f(o_j)) > 0$ (as $f(o_i) \neq f(o_j)$), there exists a threshold $\alpha_0 = \frac{m \cdot C_{ij}}{d \cdot \rho^2(f(o_i), f(o_j))}$ such that for all $\alpha > \alpha_0$, the derivative is positive, meaning $\rho(v'_i, v'_j)$ increases as $\alpha$ increases.

For a lower bound, we take the minimum value:
\begin{equation}
\rho(v'_i, v'j) \geq \frac{1}{\beta}\sqrt{\rho^2(v(o_i), v(o_j)) + \alpha^2 \cdot \frac{d}{m} \cdot \rho^2(f(o_i), f(o_j)) - 2\alpha \cdot C{ij}}
\end{equation}

\textbf{Part 4: Attribute Distance Order Preservation.}
Consider records $o_i, o_j, o_k, o_l$ with identical content vectors but different attributes. Let $v(o_i) = v(o_j) = v(o_k) = v(o_l) = v^*$, but $f(o_i) \neq f(o_j)$ and $f(o_k) \neq f(o_l)$.

For the pair $o_i, o_j$, the transformed vectors are:
\begin{align}
v'_i &= \Psi(v^, f(o_i), \alpha, \beta) = \left[\frac{v^{(1)} - \alpha f(o_i)}{\beta}, \ldots, \frac{v^{(d/m)} - \alpha f(o_i)}{\beta}\right] \\
v'_j &= \Psi(v^, f(o_j), \alpha, \beta) = \left[\frac{v^{(1)} - \alpha f(o_j)}{\beta}, \ldots, \frac{v^{(d/m)} - \alpha f(o_j)}{\beta}\right]
\end{align}

The squared distance is:
\begin{align}
\rho^2(v'i, v'j) &= \sum_{l=1}^{d/m} \sum_{h=1}^{m} \left(\frac{v^{(l)}[h] - \alpha f(o_i)[h]}{\beta} - \frac{v^{(l)}[h] - \alpha f(o_j)[h]}{\beta}\right)^2 \\
&= \sum_{l=1}^{d/m} \sum_{h=1}^{m} \left(\frac{-\alpha(f(o_i)[h] - f(o_j)[h])}{\beta}\right)^2 \\
&= \frac{\alpha^2}{\beta^2}\sum_{l=1}^{d/m} \sum_{h=1}^{m} (f(o_i)[h] - f(o_j)[h])^2 \\
&= \frac{\alpha^2}{\beta^2} \cdot \frac{d}{m} \cdot \rho^2(f(o_i), f(o_j))
\end{align}

Similarly, for the pair $o_k, o_l$:
\begin{equation}
\rho^2(v'_k, v'_l) = \frac{\alpha^2}{\beta^2} \cdot \frac{d}{m} \cdot \rho^2(f(o_k), f(o_l))
\end{equation}

Now, if $\rho(f(o_i), f(o_j)) < \rho(f(o_k), f(o_l))$, then:
\begin{equation}
\rho^2(v'_i, v'_j) = \frac{\alpha^2}{\beta^2} \cdot \frac{d}{m} \cdot \rho^2(f(o_i), f(o_j)) < \frac{\alpha^2}{\beta^2} \cdot \frac{d}{m} \cdot \rho^2(f(o_k), f(o_l)) = \rho^2(v'_k, v'_l)
\end{equation}

Taking the square root of both sides:
\begin{equation}
\rho(v'_i, v'_j) < \rho(v'_k, v'_l)
\end{equation}

This proves that the transformation $\Psi$ preserves the order of attribute distances when content vectors are identical.
\end{proof}

\end{document}
